\section{A checkpoint construction for joint guarantees}
\label{sec:method}

The upper bounds use a fixed-precision construction, Checkpoint-UCB. Its
reporting intervals are formed once, while its reward estimates use every
observation. Index arms by arrival order $i=1,2,\ldots$, with a fixed order
within each batch. Choose $0<\rho<1$ before observing rewards and set
\begin{equation}
 q_i(\rho)=\left\lceil\frac2{\rho^2}
                  \log\frac{4i(i+1)}\delta\right\rceil.
 \label{eq:checkpoint}
\end{equation}
Before this lifetime count is reached, arm $i$ has reporting interval $[0,1]$.
At its $q_i$th observation the interval becomes permanently
\begin{equation}
 [\ell_i,u_i]=[\hat r_i(q_i)-\rho/2,
                 \hat r_i(q_i)+\rho/2]\cap[0,1].
 \label{eq:fixed_interval}
\end{equation}
Although its completion time is adaptive, the
interval uses a predetermined potential prefix. Different precisions
$\rho_i$ may also be prescribed before the corresponding arms are observed;
data-dependent checkpoint selection is not covered.

At even times the policy serves the first incomplete checkpoint in arrival
order, if one remains. At every other time it maximizes, over all
$i\in\mathcal A_t$,
\begin{equation}
 J_i(t)=\hat r_i(N_i(t-1))+
             \sqrt{\frac{3\log(2t)}{N_i(t-1)}}.
 \label{eq:ucb}
\end{equation}
A zero count gives index $+\infty$; ties are broken by the smallest
arrival-order label. Both kinds of pull update the same count and reward
sum. After observing the selected reward, define
$\underline r^*(t)=\max_{i\in\mathcal A_t}\ell_i$ and
$\overline r^*(t)=\max_{i\in\mathcal A_t}u_i$, and report
\begin{equation}
 \mathcal P_t=\{i:\ell_i\ge\overline r^*(t)-\varepsilon\},\qquad
 \mathcal N_t=\{i:u_i<\underline r^*(t)-\varepsilon\},\qquad
 \mathcal V_t=\mathcal A_t\setminus(\mathcal P_t\cup\mathcal N_t).
 \label{eq:reports}
\end{equation}

\begin{algorithm}[p]
\caption{Checkpoint-UCB}
\label{alg:checkpoint-ucb}
\begin{algorithmic}[1]
\Require $\varepsilon\in(0,1)$, $\delta\in(0,1/2)$, predetermined $\rho\in(0,1)$
\State $\mathcal A^{\rm cur}\gets\varnothing$, FIFO queue $\mathcal Q^{\rm srv}\gets\varnothing$, $k\gets0$
\For{physical time $t=1,2,\ldots$}
  \For{each arm arriving at $t$, in a fixed within-batch order}
    \State $k\gets k+1$; assign label $i\gets k$; $\mathcal A^{\rm cur}\gets\mathcal A^{\rm cur}\cup\{i\}$
    \State $n_i\gets0$, $s_i\gets0$, $[\ell_i,u_i]\gets[0,1]$
    \State $q_i\gets\left\lceil(2/\rho^2)\log\{4i(i+1)/\delta\}\right\rceil$; append $i$ to $\mathcal Q^{\rm srv}$
  \EndFor
  \While{$\mathcal Q^{\rm srv}\ne\varnothing$ and $n_{\operatorname{front}(\mathcal Q^{\rm srv})}\ge q_{\operatorname{front}(\mathcal Q^{\rm srv})}$}
    \State Remove the front of $\mathcal Q^{\rm srv}$
  \EndWhile
  \If{$t$ is even and $\mathcal Q^{\rm srv}\ne\varnothing$}
    \State $I\gets\operatorname{front}(\mathcal Q^{\rm srv})$ \Comment{designated service pull}
  \Else
    \State $J_i\gets+\infty$ if $n_i=0$, and $J_i\gets s_i/n_i+\sqrt{3\log(2t)/n_i}$ otherwise, for $i\in\mathcal A^{\rm cur}$
    \State $I\gets\min\operatorname*{arg\,max}_{i\in\mathcal A^{\rm cur}}J_i$
  \EndIf
  \State Pull $I$ once and observe $x$; $n_I\gets n_I+1$, $s_I\gets s_I+x$
  \If{$n_I=q_I$}
    \State $m\gets s_I/q_I$; $\ell_I\gets\max\{0,m-\rho/2\}$, $u_I\gets\min\{1,m+\rho/2\}$
  \EndIf
  \State $L\gets\max_{i\in\mathcal A^{\rm cur}}\ell_i$, $U\gets\max_{i\in\mathcal A^{\rm cur}}u_i$
  \State $\mathcal P_t\gets\{i\in\mathcal A^{\rm cur}:\ell_i\ge U-\varepsilon\}$, $\mathcal N_t\gets\{i\in\mathcal A^{\rm cur}:u_i<L-\varepsilon\}$
  \State $\mathcal V_t\gets\mathcal A^{\rm cur}\setminus(\mathcal P_t\cup\mathcal N_t)$; report $(\mathcal P_t,\mathcal N_t,\mathcal V_t)$
\EndFor
\end{algorithmic}
\end{algorithm}

Algorithm~\ref{alg:checkpoint-ucb} observes exactly one reward per physical
round. Its variables $n_i$ and $s_i$ are the lifetime count and reward sum.
An ordinary UCB pull can complete any checkpoint, including one behind the
queue's front. The queue discards such completed entries when they reach its
front. An interval is computed only when $n_i=q_i$, and is not recomputed
when $n_i>q_i$. Reaching $q_i$ forms fixed-precision evidence; it need not
produce a determinate report if other arms are still uncertain or the arm
lies closer to the boundary than that evidence can resolve.

All arrived arms remain eligible for \eqref{eq:ucb}, including negatively
reported arms. The intervals are nonempty on every path, so
$\mathcal P_t$ and $\mathcal N_t$ are disjoint. A count, sum, and interval
are stored for every arm; a direct index or report update costs $O(K_t)$
operations. The policy uses no terminal horizon, future arrivals, true means,
or reference identities.

\begin{proposition}[Simultaneous reporting and unconditional reward cost]
\label{prop:reporting}
On the full bounded model, Checkpoint-UCB satisfies
$\Prob(E_{\rm lab})\ge1-\delta/2$. If $\mathcal I_T$ is its set of designated
service time steps, then at every deterministic $T$,
\begin{equation}
 \E R_T\le\min\left\{T,K_T+4\sqrt{3K_TT\log(2T)}+1+
                \E\sum_{t\in\mathcal I_T}\Delta_{I_t}(t)\right\}.
 \label{eq:reward_bound}
\end{equation}
Pathwise, $|\mathcal I_T|\le\sum_{i\le K_T}q_i$.
\end{proposition}

Hoeffding's inequality bounds the failure at checkpoint $i$ by
$\delta/\{2i(i+1)\}$. Summing over potential arm indices gives one coverage
event of probability at least $1-\delta/2$, on which every report is correct.
For \eqref{eq:reward_bound}, a separate, summable concentration argument at
physical time $t$ controls ordinary UCB pulls. Each arm and pre-pull count
is charged at most once; inserted service observations only skip counts.
This expectation bound holds without conditioning on reporting success,
so it omits no $\delta T$ term. It bounds one policy's total regret, rather
than a regret difference between its trajectory and a hypothetical policy
without service. \suppproof

\section{Fundamental costs of continual maintenance}
\label{sec:maintenance}

\subsection{One uncertain arrival can obstruct many reports}

The first lower bound retains one reward environment throughout the time axis.
It therefore measures a cumulative burden without adding the worst cases of
incompatible experiments. Fix
$0<\varepsilon\le1/16$, $0<\delta\le1/16$, and $0<\eta<\varepsilon/2$.
An initial reference has Bernoulli mean $1/4+\varepsilon/2$. Candidates arrive
on any prescribed schedule. In environment zero each has mean $p=1/4$;
environment $j$ changes only candidate $j$ to mean $p+2\varepsilon$.
All candidates are good in environment zero, with boundary distance
$\varepsilon/2$. After candidate $j$ arrives in environment $j$, every other
candidate is bad.

Let $M_t$ count arrived candidates and $a_t=M_t-M_{t-1}$. Define
\begin{align}
 q_-&=\left\lfloor\frac{\kl(1/4,\delta)}
                              {\kl(1/4,1/4+2\varepsilon)}\right\rfloor,
 &Q_t&=(Q_{t-1}+q_-a_t-1)_+,\quad Q_0=0,\notag\\
 \mathcal B_T&=\sum_{t\le T}(M_t-1)_+\ind\{Q_t>0\}.
 \label{eq:broadcast_queue}
\end{align}
This analytical queue is an information-capacity certificate: arrivals add
$q_-$ units per candidate, while only one observation is possible per round.
It is distinct from the algorithm's FIFO queue $\mathcal Q^{\rm srv}$ and
does not impose a pathwise sampling quota on every candidate. In the stated
range $q_-\ge2$.

\begin{theorem}[Broadcast uncertainty]
\label{thm:broadcast}
For every policy reliable on environment zero and all its one-candidate
alternatives,
\[
 \E_0\mathcal U_T(\eta)\ge(3/4-\delta)\mathcal B_T.
\]
If $\mathcal B_T>0$, then
\[
 \Prob_0\{\mathcal U_T(\eta)\ge\mathcal B_T/4\}\ge1/4.
\]
These conclusions still hold when the policy knows the full arrival schedule,
the horizon, and the reference distribution.
\end{theorem}

At a busy time, some deterministic candidate label has expected sample count
below $q_-$. Changing that candidate alone has too little history KL to permit
positive certification of another candidate with probability at least $1/4$.
Reliability also bounds incorrect negative reports. The remaining probability
must be unresolved. The same reference environment supports this argument at
every time; the difficult label is selected from expected counts rather than
from a favorable realized history. The supplementary proof includes the
bounded-variable step needed for the probability statement.

If $m$ candidates have arrived by time $s$, the queue obeys the deterministic
inequality
\begin{equation}
 \sum_{t\le s+qm-1}(M_t-1)_+\ind\{Q_t(q)>0\}
                  \ge\frac{q-1}{2}m(m-1),\qquad q\ge2.
 \label{eq:queue_area}
\end{equation}
The factor $m^2$ reflects the number of affected reports as well as the work
needed to exclude high-quality alternatives. The subtraction of one unit in
$q-1$ accounts for the last observation of each busy period: reports are issued
after service.

\subsection{Worst-schedule maintenance area}

Let $\mathfrak A_{T,K}$ be the smallest $a$ attainable by a full-model reliable
policy with
\[
 \sup_{\boldsymbol\tau:K_T\le K}\sup_{\boldsymbol r}
 \Prob_{\boldsymbol r}\{E_{\rm lab}^c\ \text{or}\ \mathcal U_T(\eta)>a\}\le\delta.
\]
This optimization imposes no regret budget. Policies may know $T,K$ and the
schedule for the lower bound. The upper bound uses neither $T$ nor $K$ as
input.

\begin{theorem}[Worst-schedule area and joint attainability]
\label{thm:worst_area}
For integers $K\ge2$ and $T\ge1$, $0<\varepsilon\le1/8$,
$0<\eta\le\varepsilon/32$, and $0<\delta\le1/16$,
\begin{equation}
 2^{-24}\min\left\{KT,\frac{K^2}{\eta^2}\log\frac K\delta\right\}
 \le\mathfrak A_{T,K}\le
 256\min\left\{KT,\frac{K^2}{\eta^2}\log\frac K\delta\right\}.
 \label{eq:worst_area}
\end{equation}
Checkpoint-UCB with $\rho=\eta/4$ attains the upper order and also satisfies
\begin{equation}
 \E R_T\le\min\left\{T,K+4\sqrt{3KT\log(2T)}+1+
                         99K\eta^{-2}\log(K/\delta)\right\}.
 \label{eq:joint_general}
\end{equation}
\end{theorem}

For the upper bound, completed checkpoints imply endpoint errors at most
$\eta/4$, sufficient to classify all arms outside the boundary band. The
unfinished-work queue is dominated by adding $q_i$ at each arrival and serving
at even rounds. On the simultaneous coverage event,
\begin{equation}
 \mathcal U_T(\eta)\le\min\left\{K_TT,
                         2K_T\sum_{i\le K_T}q_i\right\}.
 \label{eq:queue_upper}
\end{equation}
Additional UCB observations reduce the work; they are not charged again as
new checkpoint samples. The lower bound uses $K-1$ candidates with independent
means $1/2\pm2\eta$ and a reference of mean $1/2+\varepsilon$, all available
at the first pull. Simultaneous partial recovery and a horizon truncation give
the probability statement required by the definition of $\mathfrak A_{T,K}$.
The statistical $\log K$ factor has fixed-arm multiple-testing precedents;
this worst-schedule lower construction itself requires no repeated arrivals.

\subsection{Growth rates and the limits of total-count criteria}

\begin{corollary}[Arrival-growth boundary]
\label{cor:growth}
Fix the task parameters of Theorem~\ref{thm:broadcast} and one infinite arrival
schedule. If a policy is reliable for every reward environment and has a
deterministic envelope $a_T=o(T)$ satisfying \eqref{eq:joint_protocol} in
every environment, then $K_t=o(\sqrt t)$. Conversely,
$K_T^2\log(eK_T)=o(T)$ suffices for one horizon-free policy to have both
$\E R_T=o(T)$ and such a sublinear area envelope. In particular, for
$K_t\asymp t^\alpha$, $0<\alpha<1$, simultaneous attainment is possible if
and only if $\alpha<1/2$.
\end{corollary}

The necessity follows from \eqref{eq:queue_area}. For sufficiency, the general
checkpoint completion bound in Supplementary Lemma~\ref{app:completion_general}
applies throughout $0<\eta<\varepsilon/2$, together with
Proposition~\ref{prop:reporting}. This argument does not require the narrower
range $\eta\le\varepsilon/32$ in Theorem~\ref{thm:worst_area}. The exponent
statement leaves a logarithmic gap for general critical schedules.

Within the parameter range of Theorem~\ref{thm:worst_area},
\eqref{eq:worst_area} also gives a uniform terminal-count statement: over all
schedules with $K_T\le k(T)$, sublinear expected regret and sublinear service
envelopes are jointly attainable if and only if
$k(T)^2\log k(T)=o(T)$. Here the worst schedule can change with $T$.
This statement concerns a different schedule class from the single infinite
schedule in Corollary~\ref{cor:growth}.

The necessary condition $K_t=o(\sqrt t)$ is not sufficient on every fixed
schedule. For example, batches of size
\[
 m_n=\left\lfloor\frac{\sqrt{s_n}}{(\log s_n)^{1/4}}\right\rfloor,
 \qquad s_n=2^{2^n},
\]
give $K_t=o(\sqrt t)$ but preclude any $o(T)$ area envelope uniform over
reward environments. Near $s_n$, simultaneous classification of that burst
requires accumulated unresolved area of order $m_n^2\log m_n$, which exceeds
$s_n$ by a growing factor. The supplementary argument selects a single
infinite schedule; the difficult reward environment is allowed to vary with
the endpoint, as required by the stated uniform guarantee.

\section{Temporary membership tasks and regret trade-offs}
\label{sec:tradeoff}

\subsection{A repeated-update family}

We keep the reward environment, arrival schedule, endpoint, and reporting
parameters fixed when comparing two area budgets. Choose integers $J\ge1$,
$M\ge\max\{2,J\}$ and
\begin{equation}
 0<\varepsilon\le1/8,\quad 0<\delta\le1/16,\quad
 0<v,\quad Jv\le\varepsilon/4,\quad 0<d\le v/32,\quad\eta=d/2.
 \label{eq:family_parameters}
\end{equation}
At $a_e=1+eH$, $e=0,\ldots,J$, a reference with mean $b_e=1/2+ev$
arrives. For each $e<J$, the same batch contains $M$ candidates with means
\begin{equation}
 r_{e,k}=b_e-\varepsilon+\sigma_{e,k}d,
           \qquad \sigma_{e,k}\in\{-1,+1\}.
 \label{eq:family_means}
\end{equation}
All rewards are Bernoulli and all arms remain available. The last batch
contains only its reference, and $T=(J+1)H$. The family $\Theta$ includes
every sign configuration and every within-batch permutation. The policy is
not told which arm is the reference. Lower bounds may disclose those
identities and the schedule.

In its birth window a candidate is good exactly when its sign is positive.
Its reward gap is $\varepsilon-\sigma_{e,k}d$, while its distance to the
membership boundary is $d$. At the next update its gap becomes
$\varepsilon+v-\sigma_{e,k}d$, so both signs are bad, with boundary distance
at least $v-d$. Every reference remains good: the final good set consists of
all references. The tasks that distinguish the candidate signs expire even
though their arms remain available. The restriction $Jv\le\varepsilon/4$
limits the number of updates at any fixed positive improvement size.

\subsection{Low area forces enough temporary tasks to finish}

The lower-bound ingredient permits arbitrary interleaving and stopping. For
$m$ independent uniform signs with observation laws $P_i^+,P_i^-$, put
\[
 \kappa=\max_i\max\{D(P_i^+\Vert P_i^-),D(P_i^-\Vert P_i^+)\}
                                                    \in(0,\infty).
\]
If a bounded stopping time has at least half the signs determined with prior
probability at least $1/2$, and all configurations have familywise report
error at most $\delta\le1/16$, then
\begin{equation}
 \overline\E\sum_{i=1}^mN_i\ge
                 \frac{m}{1024\kappa}\log\frac m\delta.
 \label{eq:partial_recovery}
\end{equation}
Other observations must not encode these signs. The proof factors the
conditional posterior after the adaptive sampling and bounded stopping
kernels cancel, and combines simultaneous recovery with coordinatewise
change of measure \citep{kaufmann2016}. It does not assume independence of
the random sample counts. The product prior is a proof device; the final
lower bound concerns a fixed environment.

A window that never reaches half-completion contributes at least $MH/2$
area. Thus, on the event $\mathcal U_T(\eta)\le V/8$, where $V=JMH$, at
least $3J/4$ windows half-complete. A common product prior over all windows
turns this area requirement into enough instances of
\eqref{eq:partial_recovery}. Each candidate observation costs at least
$\varepsilon-d$ regret, and observations from distinct batches are charged
once. This argument does not impose a deadline on every individual batch;
the policy can choose where to defer. The heterogeneous completion
probabilities and their information-allocation relaxation are given in
Supplementary Section~\ref{app:area_lower}. That relaxation is a necessary
constraint, with no general causal attainability claim.

\subsection{Shared evidence for area and reward guarantees}

Checkpoint-UCB may use a predetermined precision $\rho_e$ for each arriving
batch. Consider the heterogeneous repeated-update family: batch $F_e$ has
reference mean $b_e$ and $M_e$ candidate means
$b_e-\varepsilon\pm d_e$, all in the allowed reward range. The windows
partition physical times $1,\ldots,T$; write $a_e$ for the arrival time of
batch $e$, $H_e$ for its window length, and
$S_e=\sum_{i\in F_e}q_i(\rho_e)$ for its total checkpoint work; the final
batch contains only a reference. Set $M_{-1}=M_J=0$. Suppose
\begin{equation}
 H_e\ge2S_e,\qquad b_J-b_0\le\varepsilon/4,\qquad
 d_e+\rho_e+\rho_{e+1}<b_{e+1}-b_e,\qquad
 2\max_e\rho_e\le\varepsilon/2.
 \label{eq:reuse_conditions}
\end{equation}
The neighboring-batch condition is imposed for $e<J$; the window and
precision bounds include the final batch. These are conditions on the
performance family, not future information supplied to the policy. Put
$\mathcal F=\{e<J:2\rho_e<d_e\}$. For any arm that has a positive gap at
some time through $T$, define
\[
 \gamma_i=\min_{t\le T:i\in\mathcal A_t,\,\Delta_i(t)>0}\Delta_i(t),
 \qquad
 D_i^T=\max_{t\le T:i\in\mathcal A_t}\Delta_i(t).
\]
Only such arms are included in the sum containing $\gamma_i$ below.

\begin{proposition}[Shared update from lifetime evidence]
\label{thm:reuse}
Under \eqref{eq:reuse_conditions}, on the checkpoint coverage event,
\begin{equation}
 \mathcal U_T(\eta)\le
 \sum_{e\notin\mathcal F}M_eH_e+
 2\sum_{e\in\mathcal F}M_eS_e+
 2\sum_{e=0}^J(e+1+M_{e-1})S_e.
 \label{eq:reuse_area}
\end{equation}
When the next batch completes its checkpoints, every preceding-batch candidate
can be excluded using its saved interval and the new reference lower bound,
without another observation of that old candidate for certification.
For the same policy, the unconditional expected regret satisfies
\begin{equation}
 \begin{aligned}
 \E R_T\le{}&\sum_{e<J}\sum_{k=1}^{M_e}(\varepsilon+d_e)q_{e,k}\\
 &+\sum_{i:\gamma_i\text{ exists}}D_i^T
 \left[\left\lceil\frac{12\log(2T)}{\gamma_i^2}\right\rceil-q_i+1\right]_+
 +1.
 \end{aligned}
 \label{eq:reuse_regret}
\end{equation}
\end{proposition}

The last area term charges temporarily affected references and the preceding
candidate cohort, rather than all historical candidates at every update.
For $d_e>\eta$, batches in $\mathcal F$ complete their current fine
membership task; a batch with $d_e\le\eta$ lies inside the evaluation band
in its birth window, and its corresponding area term is conservative.
An old candidate's upper endpoint and the new reference's lower endpoint
differ by more than $\varepsilon$, giving the later negative report without
new observations of that candidate. This does not prevent the reward layer
from selecting it again.

Every $q_i$ is a lifetime quota. The condition $H_e\ge2S_e$ ensures that all
observations up to that quota occur in the arm's birth window. The first
term of \eqref{eq:reuse_regret} charges candidate observations in that
window, while reference observations then have zero regret. The second
term counts excess ordinary UCB pulls using distinct pre-pull counts. The
expectation is not conditioned on checkpoint coverage. The birth-window
condition is essential to this charging and does not cover arbitrary
overlapping arrivals.

\subsection{Rapid and deferred maintenance at the same endpoint}

For the homogeneous family
\eqref{eq:family_parameters}--\eqref{eq:family_means}, set
\begin{gather}
 K=J(M+1)+1,\quad L_K=\log\frac{4K(K+1)}\delta,\quad
 \rho_f=d/4,\quad\rho_c=v/8,\notag\\
 q_f=\lceil32L_K/d^2\rceil,\quad q_c=\lceil128L_K/v^2\rceil,\quad
 H=128(M+J+1)q_f,\quad V=JMH,\quad T=(J+1)H.
 \label{eq:chosen_horizon}
\end{gather}
The quantities $q_f,q_c$ are uniform analysis bounds on the actual armwise
quotas. Neither final arm count nor future times are policy inputs. Fine
precision can be set from the task input as $\rho_f=\eta/2$; the coarse
precision is a separate preselected configuration. Automatic adaptation to
an unknown improvement scale is not asserted.

\begin{theorem}[Rapid and deferred maintenance at the same endpoint]
\label{thm:tradeoff}
On the preceding family and horizon,
\begin{equation}
 2^{-20}\frac{\varepsilon JM}{d^2}\log\frac M\delta
 \le\mathfrak R_\Theta(V/8)\le
 2^{18}\frac{\varepsilon JM}{d^2}\log\frac M\delta.
 \label{eq:rapid_matched}
\end{equation}
The coarse-precision Checkpoint-UCB policy has joint label-and-area success
probability at least $1-\delta$ with area budget $9V/8$, and
\begin{equation}
 \sup_{\theta\in\Theta}\E_\theta R_T\le
 2\varepsilon JM q_c+
 64\left(\frac{JM}{\varepsilon}+\frac{J^2}{v}\right)\log(2T)+1.
 \label{eq:coarse_regret}
\end{equation}
Consequently, at fixed legal $J,M,\varepsilon,v,\delta$ and $d\downarrow0$,
with $\eta=d/2$ and the horizon in \eqref{eq:chosen_horizon},
$\mathfrak R_\Theta(V/8)=\Theta(d^{-2})$, whereas
$\mathfrak R_\Theta(9V/8)=O(\log(1/d))$.
\end{theorem}

For the lower bound, low area forces enough windows to satisfy
\eqref{eq:partial_recovery}. For the upper bound, \eqref{eq:reuse_area}
gives the two area budgets. At the specified horizon the deterministic bound
$R_T\le2\varepsilon T$ suffices for the rapid order; it would not justify
the same conclusion at arbitrarily longer windows. The coarse bound instead
counts checkpoint observations and excess ordinary UCB pulls separately,
as in \eqref{eq:reuse_regret}.

The two budgets concern the same $\Theta,\varepsilon,\eta,\delta,T$, but
allow different service levels. Their ratio does not describe an individual
delay for every arm. The theorem does not claim that the coarse policy meets
the rapid budget, that its regret upper bound is matched, or that all
intermediate budgets are characterized. A preselected mixture of fine and
coarse batches gives additional attainable points in the supplement. The
separation concerns this specified repeated-update family, with $\eta$ and
$T$ changing together as $d\downarrow0$.

\subsection{A probability lower bound under the same area requirement}

An expectation lower bound alone could reflect unusual sample paths. The
following result gives a stronger conclusion without replacing cumulative
area by individual deadlines.

\begin{theorem}[A nonvanishing probability of the rapid-maintenance cost]
\label{thm:quantile_main}
In the repeated-update family, suppose a reliable policy satisfies
\[
 \Prob_\sigma\{E_{\rm lab},\mathcal U_T(\eta)\le V/8\}\ge1-\delta
\]
for every sign configuration. Then one fixed configuration satisfies
\begin{equation}
 \Prob_\sigma\left\{R_T\ge
 2^{-20}\frac{\varepsilon JM}{d^2}\log\frac M\delta\right\}
 \ge\frac{11}{80}.
 \label{eq:quantile_main}
\end{equation}
\end{theorem}

The proof stops each window at half-completion, its endpoint, or its
$(q+1)$st candidate observation, where $q$ is a fixed sample budget. Stopping
at the first observation beyond $q$ retains reports that become determinate
after a reference observation while the candidate count is still $q$.
Applying the partial-recovery lemma to this bounded experiment makes the
probability of a cheap completion less than $1/2$. Low total area requires
many completions; counting cheap ones then leaves enough expensive windows
with probability at least $1/5-\delta\ge11/80$. Their actual observations
give \eqref{eq:quantile_main}. No independence between window-completion
events is needed. Averaging under the common prior then selects one fixed
environment.

Since $11/80>\delta$, the rapid endpoint in \eqref{eq:rapid_matched} also
has matching order when regret and area must be bounded on one event of
probability at least $1-\delta$. The exact joint-quantile definition and
corollary are in Supplementary Section~\ref{sec:quantile_extension}. This
consequence concerns the rapid endpoint at the prescribed horizon and gives
no matching joint-quantile result for deferred maintenance.
